% !TeX program = pdflatex
% ============================================================
%  paper_extended.tex
%  Accuracy Evaluation of MediaPipe Pose Estimation
%  Using a Unity Environment --- Extended 8-page Review Paper
%
%  Base settings inherited from: docs/paper/original_paper.md
%  Target journal: IEEE Access / Sensors (MDPI)
%
%  Overleaf upload structure (upload entire source/ folder as project root):
%    main.tex                         <- compile this
%    figs/
%      fig01_camera_layout.png
%      fig02a_unity_env.png
%      fig02b_frame_sample.jpg
%      fig03_concept_metrics.png
%      fig04a_elbow_L_theta_y05.jpg
%      fig04b_elbow_R_theta_y05.jpg
%      fig05a_hip_L_psi_y05.jpg
%      fig05b_hip_R_psi_y05.jpg
%      fig06a_corr_theta.png
%      fig06b_corr_psi.png
%      fig07a_mae_elbow_L_y05.png
%      fig07b_mae_elbow_L_y15.png
%      fig07c_mae_elbow_L_y20.png
% ============================================================

\documentclass[twocolumn]{article}

% ---- Encoding (pdfLaTeX) ----------------------------------
\usepackage[utf8]{inputenc}
\usepackage[T1]{fontenc}
\usepackage{gensymb}  % \degree works in both text and math mode

% ---- Package loading (original_paper.md baseline) ----------
\usepackage{graphicx}
\usepackage{geometry}
\usepackage{amsmath}
\usepackage{booktabs}
\usepackage{float}
\usepackage{url}
\usepackage{titlesec}

% ---- Extended packages (added for 8-page version) ----------
\usepackage{subcaption}   % subfigure layout (Fig 4, Fig 5, Fig 7/8 pairs)
\usepackage{multirow}     % merged rows in tables
\usepackage{array}        % column formatting
\usepackage{siunitx}      % angle unit \ang{}
\usepackage[hidelinks]{hyperref}  % clickable refs, no colored boxes

% ---- Page geometry (inherited from original_paper.md) ------
\geometry{top=24mm,bottom=24mm,left=19mm,right=19mm}
\setlength{\parskip}{0.15em}
\setlength{\floatsep}{6pt}
\setlength{\textfloatsep}{8pt}
\titlespacing*{\section}   {0pt}{2ex plus 0.5ex minus 0.2ex}{1.2ex plus 0.2ex}
\titlespacing*{\subsection}{0pt}{1.8ex plus 0.5ex minus 0.2ex}{0.9ex plus 0.2ex}
\titlespacing*{\subsubsection}{0pt}{1.5ex plus 0.3ex minus 0.2ex}{0.6ex plus 0.2ex}

\sloppy

% ---- Title block -------------------------------------------
\title{Accuracy Evaluation of MediaPipe Pose Estimation\\
       Using a Unity Simulation Environment}

\author{%
  Non-member\quad Fumima Taira$^{*a)}$ \qquad
  Member\quad    Tsukasa Kato$^{*b)}$ \qquad
  Non-member\quad Jin Afuso$^{*c)}$\footnotemark[1]
}
\date{}

% ============================================================
\begin{document}
\maketitle

\begingroup
\renewcommand{\thefootnote}{}
\footnotetext[1]{%
  $^{*a)}$ Okinawa Prefectural Kaiko High School,\quad
  $^{*b)}$ Graduate School of Education, University of the Ryukyus,\quad
  $^{*c)}$ lollol Inc.}
\endgroup

% ============================================================
% ABSTRACT  (target: ~0.5 page)
% ============================================================
\section*{Abstract}

%%TODO%% --- Expand to 4--5 sentences. Include:
%  - Background: MediaPipe accuracy, viewpoint limitation of prior work
%  - Method: 505 cameras, 107 frames, 259,356 observations, 12 joints
%  - Two error types: Joint Angle MAE + Direction Angle Error (Δθ, Δψ)
%  - Key results (quantitative): elbow Δθ ~58\degree{}, hip Δψ r=−0.84, shoulder improvement after Y-flip
%  - Contributions as bullet list (3--4 items)

This study presents a large-scale quantitative evaluation of MediaPipe pose
estimation accuracy using Ground Truth collected in a Unity simulation
environment. Prior evaluations have largely been limited to frontal or
restricted camera configurations; evaluation across diverse viewpoints
essential in practical deployments has been insufficient.

We collected walking data from 505 multi-view cameras over 107 frames,
yielding 259,356 joint observations across 12 joints. Two error types were
evaluated: (1)~joint angle error (MAE) and (2)~direction angle error
($\Delta\theta$, $\Delta\psi$), together with inter-joint correlation
analysis. Our main findings are as follows.
\begin{itemize}\itemsep0pt
  \item Unifying the coordinate system (Y-axis orientation) reduces the mean
        elbow direction angle error from $\approx121\degree{}$ to $\approx58\degree{}$, and
        shoulder error from $\approx120\degree{}$ to $\approx10\degree{}$.
  \item Left and right hip depth errors exhibit a strong negative correlation
        ($r = -0.84$), revealing a systematic pelvic rotation artefact.
  \item Upper-limb errors are strongly coupled ($r = 0.72$--$0.77$),
        suggesting hierarchical error propagation.
  \item Error patterns are viewpoint-dependent and largely predictable,
        offering scope for post-processing correction.
\end{itemize}

\vspace{-0.8em}

% ============================================================
% 1. INTRODUCTION  (target: ~1.5 pages)
% ============================================================
\section{Introduction}

Pose estimation has become a core enabling technology across sports
analysis, clinical rehabilitation, ergonomics, and human--computer
interaction~\cite{mediapipe,blazepose}. MediaPipe, developed by Google, is
particularly attractive for practitioners because it enables real-time
processing on commodity hardware without GPU acceleration~\cite{mediapipe}.

Despite its widespread deployment, a systematic and multi-viewpoint
quantitative evaluation of MediaPipe's 3D accuracy remains lacking.
Bazarevsky~et~al.~\cite{blazepose} report BlazePose benchmark performance,
but their evaluation is restricted to a curated benchmark with limited
viewpoint diversity. In real-world deployments, camera placement is
constrained by the environment and cannot be controlled; understanding how
accuracy varies with viewpoint is therefore essential.

%%TODO%% --- Expand this section with:
%  (a) Prior work on pose estimation accuracy evaluation (see Related Work
%      section --- cross-reference here briefly)
%  (b) The monocular depth ambiguity problem
%  (c) Role of simulation/Unity-based GT collection
%  (d) Gap this paper fills

The contributions of this paper are as follows:
\begin{enumerate}\itemsep0pt
  \item \textbf{Large-scale multi-viewpoint evaluation}: 505 camera
        positions spanning four height layers and a wide horizontal range,
        yielding 259,356 observations --- the largest multi-viewpoint
        MediaPipe evaluation to date.
  \item \textbf{Coordinate system unification}: We identify and correct a
        systematic Y-axis orientation mismatch between Unity and MediaPipe,
        and demonstrate its substantial effect on elbow direction accuracy.
  \item \textbf{Direction angle error analysis}: We introduce and quantify
        $\Delta\theta$ and $\Delta\psi$ as complementary metrics to joint
        angle MAE, revealing bias patterns invisible to MAE alone.
  \item \textbf{Inter-joint error correlation}: Pearson correlation analysis
        across 12 joints exposes systematic model-level constraints including
        pelvic anti-correlation and upper-limb coupling.
\end{enumerate}

% ============================================================
% 2. RELATED WORK  (target: ~1 page) --- NEW SECTION
% ============================================================
\section{Related Work}

\subsection{MediaPipe and BlazePose Accuracy}

BlazePose, the neural-network backbone of MediaPipe Pose, was introduced by
Bazarevsky~et~al.~\cite{blazepose}, who report PCK@0.2 of 72.0\,\% on the
COCO validation set and state-of-the-art performance on a proprietary
fitness-video benchmark. Their evaluation is, however, restricted to
near-frontal perspectives typical of fitness recordings; performance under
diverse camera viewpoints is not assessed.

Independent comparative evaluations have confirmed the accuracy--speed
trade-off of MediaPipe. Mroz~et~al.~\cite{mroz2021} compared BlazePose with
OpenPose~\cite{openpose} on a standardised benchmark and found BlazePose to
be substantially faster (0.008\,s vs.\ 0.107\,s per frame on the same
hardware) at the cost of slightly lower accuracy (PCK 0.77 vs.\ 0.83).
Although these results confirm that MediaPipe Pose is well-suited for
latency-critical applications, neither the original BlazePose benchmark nor
subsequent comparative evaluations provide systematic accuracy data across
diverse camera positions. The present study addresses this gap with a
505-camera 3-D grid spanning four height layers (Y\,=\,0.5--2.0\,m).

\subsection{Simulation-Based Ground Truth Collection}

Obtaining accurate 3-D ground truth at scale is impractical in real-world
settings: manual 3-D annotation is labour-intensive and optical
motion-capture systems are costly and laboratory-confined. Synthetic
rendering has therefore become a standard approach. Varol~et~al.~\cite{surreal}
introduced SURREAL, rendering over 6~million photo-realistic frames of
SMPL-body avatars with perfect 2-D/3-D pose, depth, and segmentation labels;
models pre-trained on SURREAL transfer effectively to real-world benchmarks
including Human3.6M~\cite{h36m}. Patel~et~al.~\cite{agora} extended this
paradigm with AGORA, compositing 4,240 high-quality textured body scans into
natural outdoor environments and demonstrating that training on realistic
synthetic data improves regression performance on the 3DPW benchmark.

Game-engine-based capture systems add the further advantage of full camera
controllability. Ebadi~et~al.~\cite{peoplesanspeople} released
PeopleSansPeople, a Unity-based synthetic data generator that produces
COCO-format keypoint annotations; pre-training on PeopleSansPeople data and
fine-tuning on real COCO images achieves keypoint AP of 63.5, surpassing a
real-data-only baseline of 62.0~AP in the high-data regime. Our study
similarly exploits a Unity simulation environment to record exact per-frame
3-D joint coordinates from 505 systematically placed cameras---a viewpoint
sweep that physical motion-capture installations cannot economically
reproduce.

\subsection{Monocular 3D Pose Estimation and Depth Ambiguity}

Recovering 3-D pose from a single RGB camera is inherently ill-posed: many
3-D configurations project to the same 2-D image under perspective
projection. Mehta~et~al.~\cite{mehta2017} improved generalisation to
in-the-wild scenes through multi-dataset transfer learning and introduced the
MPI-INF-3DHP benchmark, which includes greater viewpoint diversity than
Human3.6M~\cite{h36m}; nevertheless, lateral and overhead perspectives
remain under-represented in standard training corpora.

Public datasets such as Human3.6M~\cite{h36m} and COCO~\cite{coco} are
dominated by near-frontal, standing or walking poses, embedding a strong
frontal viewpoint prior in learned models. This bias causes accuracy to
degrade for cameras positioned laterally or overhead---a pattern our
multi-height evaluation confirms quantitatively (Table~\ref{tab:mae_by_layer}).
Beyond viewpoint bias, MediaPipe's skeleton model represents the pelvis as a
simple line segment connecting landmarks~23 (left hip) and~24 (right hip),
without a rigid multi-segment structure. A small horizontal pixel offset in
the 2-D image can therefore be misinterpreted as a large pelvic rotation
during 3-D lifting, producing the symmetric anti-correlated hip depth error
($r=-0.840$) reported in Section~4.

\subsection{Skeletal Constraint Models}

Several architectures impose anatomical or temporal constraints to mitigate
depth ambiguity. Mehta~et~al.~\cite{vnect} demonstrated real-time monocular
3-D pose estimation with a kinematic skeleton fitting stage that enforces
bone-length consistency and temporal stability across frames.
PoseFormer~\cite{poseformer} is a purely transformer-based video pose
estimator that captures long-range spatio-temporal joint relations without
convolutional layers, achieving 44.3\,mm MPJPE on Human3.6M.
MotionBERT~\cite{motionbert} pre-trains a Dual-stream Spatio-temporal
Transformer to recover 3-D motion from noisy 2-D observations, embedding
geometric and physical priors that generalise to mesh recovery and action
recognition, achieving 39.2\,mm MPJPE on Human3.6M.

Despite these advances, pelvic rigidity is rarely enforced as an explicit
hard constraint: the pelvis is typically modelled as a single segment or
handled implicitly through temporal consistency. The hip anti-correlation
($r=-0.840$) observed in our data suggests that constraining left- and
right-hip depth to be mutually consistent under a rigid pelvic model could
substantially reduce depth errors for walking motions---a direction we
recommend for future model design.

% ============================================================
% 3. EXPERIMENTAL METHOD  (target: ~1.5 pages)
% ============================================================
\section{Experimental Method}

\subsection{Experimental Environment}

Using Unity (version 6000.0.60f1), we constructed a walking simulation
environment with a Y-Bot humanoid avatar. The avatar performed a straight
walking motion from coordinates $(0,\,0,\,-3)$ to $(0,\,0,\,3)$ along the
positive Z-axis.

\begin{figure}[!t]
\centering
\includegraphics[width=0.30\textwidth]{figs/fig01_camera_layout.png}
\caption{Camera layout (505 positions across four height layers Y = 0.5,
  1.0, 1.5, 2.0). The centre arrow indicates the avatar's walking direction
  (+Z axis).}
\label{fig:camera_layout}
\end{figure}

Cameras were placed on a 3-D grid covering four height layers
(Y = 0.5, 1.0, 1.5, 2.0\,m) and horizontal positions
(X, Z $\in [-6, 6]$\,m). Within each layer a $13\times13$ grid was used
with the central $5\times5$ block removed, giving a theoretical maximum of
$4(13^2-5^2)=576$ positions. Data were collected at 505 of these positions.
Each position yielded 107 captured frames at a resolution of
$1280\times720$\,pixels, for a total of 259,356 joint observations
(107 frames $\times$ 505 cameras $\times$ 12 joints, after visibility
filtering). The same GT and image data are reusable for evaluating
alternative estimators such as YOLO-Pose.

%%TODO%% --- Add 1--2 sentences on avatar parameters (Y-Bot, animation clip name,
%  walking speed/cycle if known).

% NOTE (NEW_PAPER_MANAGEMENT):
%   Images/Unity実験環境.png および Images/frame_0018.jpg は
%   元の投稿論文（docs/paper/Zeval_IEEJ_submit.pdf）に掲載済みの図。
%   Zeval_DataSet_organized/07_paper/Images/ への配置が未完了。
%   元論文から抽出するか、Unity/1_Output_Photos から取得すること。
\begin{figure}[!t]
  \centering
  \includegraphics[width=0.44\textwidth]{figs/fig02a_unity_env.png}
  \vspace{0.3em}
  \includegraphics[width=0.44\textwidth]{figs/fig02b_frame_sample.jpg}
  \caption{Unity experiment environment (top) and a representative captured
    frame (bottom).}
  \label{fig:unity_env}
\end{figure}

\subsection{Capture System}

A Trigger Zone Capture System was implemented in Unity.
\texttt{CaptureSystemManager} handled character generation, animation
playback, per-frame image capture, and joint coordinate recording.
\texttt{AutoCaptureManager} read the 505 camera positions from a CSV file
and triggered automatic capture at each position. This ensured consistent,
reproducible multi-viewpoint data collection with a single character motion
sequence.

\subsection{Data Collection}

Joint coordinates (GT) were exported per frame in CSV format using Unity's
\texttt{HumanBodyBones} API~\cite{unity}. MediaPipe's Pose Landmarker
(\texttt{model\_complexity=1}, \texttt{static\_image\_mode=True},
\texttt{min\_detection\_confidence=0.5}) was applied to each saved image to
extract 3-D coordinates for 33 landmarks. Landmarks with a visibility score
below 0.5 were excluded from analysis to limit the effect of heavy
occlusion. The 12 joints retained for analysis are: left/right shoulder,
elbow, wrist, hip, knee, and ankle.

\subsection{Angle Computation and Error Evaluation}

We report \textbf{two complementary error metrics}.

\subsubsection{Joint Angle Error (Joint Angle MAE)}

The flexion angle at each joint was computed from three anatomical points:
upper arm--elbow--forearm for the elbow; thigh--knee--lower leg for the knee;
and equivalent triplets for shoulder and hip. The angle $\alpha$ at the
vertex between the two vectors was computed as
\[
  \alpha = \arccos\!\left(\frac{\mathbf{u}\cdot\mathbf{v}}
                               {\|\mathbf{u}\|\,\|\mathbf{v}\|}\right),
  \quad \alpha \in [0\degree{},\,180\degree{}].
\]
The mean absolute difference between GT and MediaPipe angles over all valid
frames and cameras gives the MAE for each joint.

\subsubsection{Direction Angle Error ($\Delta\theta$, $\Delta\psi$)}

Following the BlazePose output specification~\cite{blazepose}, relative
joint positions are expressed with the hip midpoint as origin. In this
hip-centred coordinate system we define
\[
  \theta = \mathrm{arctan2}(y,\,x), \quad
  \psi   = \mathrm{arctan2}(z,\,x),
\]
where $(x,y,z)$ is the joint position relative to the hip centre.
The angular errors $\Delta\theta$ and $\Delta\psi$ are the signed
differences between GT and MediaPipe values, normalised to
$(-180\degree{}, +180\degree{}]$.

\subsubsection{Coordinate System Unification}

Unity uses a left-handed coordinate system with Y-axis pointing
upward, whereas MediaPipe's \texttt{pose\_landmarks} uses an image
coordinate system with Y-axis pointing downward. This discrepancy
introduces a systematic sign flip in $\theta = \mathrm{arctan2}(y,x)$.
Before computing direction angle errors we applied the transformation
\[
  y_{\text{MediaPipe}}^{\text{unified}} = -y_{\text{MediaPipe}},
\]
implemented in
\texttt{06\_direction\_detection/scripts/coordinate\_transform.py}.
All direction angle results and inter-joint correlation analyses reported
below use coordinates after this unification.

\begin{figure}[!t]
\centering
\includegraphics[width=0.48\linewidth]{figs/fig03_concept_metrics.png}
\caption{Concept diagram of the two evaluation metrics. Left: joint angle
  (MAE) --- flexion angle of three points ($0$--$180\degree{}$). Right: direction
  angle ($\Delta\theta$, $\Delta\psi$) --- angular offset relative to the
  hip centre ($-180\degree{}$ to $+180\degree{}$). The figure is a simplified schematic;
  the actual analysis uses all 33 MediaPipe landmarks~\cite{blazepose}.}
\label{fig:concept_metrics}
\end{figure}

% ============================================================
% 4. RESULTS  (target: ~2.5 pages)
% ============================================================
\section{Results}

\subsection{Joint Angle MAE}

Table~\ref{tab:joint_angle_error} summarises MAE, median, and maximum
error for eight major joints. Elbow and knee joints exhibit the lowest
errors (~18\degree{} and ~17--19\degree{}, respectively), while the shoulder joints show
the highest MAE (~39--40\degree{}), likely reflecting the increased depth ambiguity
of proximal joints relative to the camera.

\begin{table}[!t]
\centering
\caption{Statistics of joint angle error (Joint Angle MAE, degrees).
  Source: \texttt{04\_mae\_heatmap/coordinate\_angle\_mae.csv}.}
\label{tab:joint_angle_error}
\begin{tabular}{lccc}
\toprule
Joint          & MAE    & Median & Max    \\
               & (deg)  & (deg)  & (deg)  \\
\midrule
Left shoulder  & 39.8   & 39.2   & 70.1   \\
Right shoulder & 39.0   & 40.0   & 58.2   \\
Left elbow     & 18.6   & 15.9   & 56.9   \\
Right elbow    & 18.2   & 13.6   & 46.2   \\
Left hip       & 30.2   & 29.4   & 66.0   \\
Right hip      & 33.3   & 32.2   & 81.7   \\
Left knee      & 17.0   & 16.6   & 43.4   \\
Right knee     & 19.5   & 19.4   & 62.3   \\
\bottomrule
\end{tabular}
\end{table}

%%TODO%% --- Add Y-layer comparison sub-table or figure:
%  Y=0.5,1.5 vs Y=1.0,2.0 split of MAE per joint.
%  Source: 04_mae_heatmap/Y=0.5,1.5/ and Y=1.0.2.0/ outputs.
%  Data to be filled in by data-collection agent --- see DATA_COLLECTION_AGENT_INSTRUCTIONS.md §2.

\subsection{Direction Angle Error ($\Delta\theta$, $\Delta\psi$)}

Table~\ref{tab:direction_angle_error} presents the direction angle error
statistics for all 12 joints after coordinate system unification.

\begin{table}[!t]
\centering
\caption{Direction angle error statistics after coordinate unification
  (degrees; absolute-value means). Source:
  \texttt{11\_direction\_ditection/output/processed\_data/joint\_summary.csv}.
  $^\dagger$Hip $|\Delta\psi|$ from separate analysis (see text).}
\label{tab:direction_angle_error}
\small
\begin{tabular}{lcccc}
\toprule
Joint & \multicolumn{2}{c}{$|\Delta\theta|$} & \multicolumn{2}{c}{$|\Delta\psi|$} \\
\cmidrule(lr){2-3}\cmidrule(lr){4-5}
      & Mean & SD & Mean & SD \\
\midrule
Left shoulder  & 10.4 & 11.7 & 94.1 & 86.4 \\
Right shoulder &  9.3 & 10.1 & 92.7 & 92.0 \\
Left elbow     & 58.3 & 24.1 & 87.3 & 90.6 \\
Right elbow    & 59.3 & 22.7 & 88.6 & 94.5 \\
Left wrist     & 84.3 & 103.7 & 88.6 & 91.9 \\
Right wrist    & 91.5 & 111.9 & 91.6 & 99.2 \\
Left hip$^\dagger$ & --- & --- & 88.9 & --- \\
Right hip$^\dagger$& --- & --- & 91.1 & --- \\
Left knee      & 17.4 & 18.8 & 84.7 & 97.9 \\
Right knee     & 18.5 & 19.1 & 87.2 & 97.5 \\
Left ankle     & 11.3 & 14.0 & 90.6 & 105.8 \\
Right ankle    & 11.9 & 14.5 & 100.4 & 103.1 \\
\bottomrule
\end{tabular}
\end{table}

Figure~\ref{fig:elbow_heatmaps} shows the camera-position-wise distribution
of elbow $\Delta\theta$.

\begin{figure}[!t]
  \centering
  \begin{subfigure}[b]{0.48\linewidth}
    \includegraphics[width=\linewidth]%
      {figs/fig04a_elbow_L_theta_y05.jpg}
    \caption{Left elbow $\Delta\theta$ (Y=0.5)}
  \end{subfigure}\hfill
  \begin{subfigure}[b]{0.48\linewidth}
    \includegraphics[width=\linewidth]%
      {figs/fig04b_elbow_R_theta_y05.jpg}
    \caption{Right elbow $\Delta\theta$ (Y=0.5)}
  \end{subfigure}
  \caption{Heatmaps of elbow direction angle error $\Delta\theta$ (Y=0.5).
    Left and right elbows show opposite-sign patterns, consistent with the
    strong negative inter-joint correlation ($r=-0.862$).}
  \label{fig:elbow_heatmaps}
\end{figure}

Figure~\ref{fig:hip_psi_heatmap} shows the hip $\Delta\psi$ heatmap.
Left and right hip exhibit near-identical absolute-value patterns due to
their strong negative correlation ($r=-0.84$).

\begin{figure}[!t]
  \centering
  \begin{subfigure}[b]{0.48\linewidth}
    \includegraphics[width=\linewidth]{figs/fig05a_hip_L_psi_y05.jpg}
    \caption{Left hip $\Delta\psi$ (Y=0.5)}
  \end{subfigure}\hfill
  \begin{subfigure}[b]{0.48\linewidth}
    \includegraphics[width=\linewidth]{figs/fig05b_hip_R_psi_y05.jpg}
    \caption{Right hip $\Delta\psi$ (Y=0.5)}
  \end{subfigure}
  \caption{Heatmaps of hip direction angle error $\Delta\psi$ (Y=0.5).
    Left and right hips show near-identical absolute-value patterns
    because $\Delta\psi_{\text{L}} \approx -\Delta\psi_{\text{R}}$
    ($r = -0.840$): when one hip is forward, the other is backward.}
  \label{fig:hip_psi_heatmap}
\end{figure}

Table~\ref{tab:mae_by_layer} breaks down joint angle MAE by camera height
layer. Errors are generally lowest at Y=1.0 and increase at Y=2.0
(overhead view), suggesting that top-down perspectives reduce estimation
accuracy, consistent with underrepresentation of such viewpoints in training
data. Figure~\ref{fig:mae_layer_compare} visualises this trend for the left
elbow across three height layers.

\begin{figure}[!t]
  \centering
  \begin{subfigure}[b]{0.32\linewidth}
    \includegraphics[width=\linewidth]{figs/fig07a_mae_elbow_L_y05.png}
    \caption{Y=0.5}
  \end{subfigure}\hfill
  \begin{subfigure}[b]{0.32\linewidth}
    \includegraphics[width=\linewidth]{figs/fig07b_mae_elbow_L_y15.png}
    \caption{Y=1.5}
  \end{subfigure}\hfill
  \begin{subfigure}[b]{0.32\linewidth}
    \includegraphics[width=\linewidth]{figs/fig07c_mae_elbow_L_y20.png}
    \caption{Y=2.0}
  \end{subfigure}
  \caption{Joint angle MAE heatmaps for the left elbow at three camera
    height layers. Error increases at Y=2.0 (overhead), consistent with
    training-data viewpoint bias. Source:
    \texttt{4\_MAE\_HEATMAP/}.}
  \label{fig:mae_layer_compare}
\end{figure}

\begin{table}[!t]
\centering
\caption{Joint angle MAE (deg) by camera height layer.
  Source: \texttt{4\_MAE\_HEATMAP/Y=0.5,1.5/} and
  \texttt{4\_MAE\_HEATMAP/Y=1.0.2.0/coordinate\_angle\_mae.csv}.}
\label{tab:mae_by_layer}
\small
\begin{tabular}{lcccc}
\toprule
Joint & Y=0.5 & Y=1.0 & Y=1.5 & Y=2.0 \\
\midrule
Left shoulder  & 39.8 & 38.9 & 39.6 & 41.0 \\
Right shoulder & 39.3 & 38.1 & 38.6 & 40.2 \\
Left elbow     & 18.1 & 17.7 & 18.3 & 20.2 \\
Right elbow    & 17.5 & 16.7 & 18.2 & 20.2 \\
Left hip       & 29.6 & 29.0 & 29.9 & 32.5 \\
Right hip      & 32.0 & 31.7 & 32.8 & 36.7 \\
Left knee      & 16.1 & 16.6 & 17.7 & 17.8 \\
Right knee     & 18.9 & 18.7 & 19.6 & 20.6 \\
\bottomrule
\end{tabular}
\end{table}

\subsection{Inter-joint Error Correlation}

Tables~\ref{tab:correlation_theta} and \ref{tab:correlation_psi} list
high-correlation pairs ($|r|>0.7$) for $\Delta\theta$ and $\Delta\psi$,
respectively.

\begin{table}[!t]
\centering
\caption{High-correlation joint pairs in XY plane ($\Delta\theta$).}
\label{tab:correlation_theta}
\small
\begin{tabular}{llr}
\toprule
Joint 1        & Joint 2        & Correlation \\
\midrule
LEFT\_ELBOW    & RIGHT\_ELBOW   & $-0.862$ \\
LEFT\_ELBOW    & LEFT\_SHOULDER &  $0.788$ \\
RIGHT\_ANKLE   & RIGHT\_KNEE    &  $0.767$ \\
LEFT\_ANKLE    & LEFT\_KNEE     &  $0.766$ \\
RIGHT\_ELBOW   & RIGHT\_SHOULDER&  $0.710$ \\
\bottomrule
\end{tabular}
\end{table}

\begin{table}[!t]
\centering
\caption{High-correlation joint pairs in XZ plane ($\Delta\psi$).}
\label{tab:correlation_psi}
\small
\begin{tabular}{llr}
\toprule
Joint 1        & Joint 2         & Correlation \\
\midrule
LEFT\_HIP      & RIGHT\_HIP      & $-0.840$ \\
RIGHT\_ELBOW   & RIGHT\_SHOULDER &  $0.770$ \\
RIGHT\_ELBOW   & RIGHT\_WRIST    &  $0.769$ \\
LEFT\_ELBOW    & LEFT\_SHOULDER  &  $0.768$ \\
RIGHT\_SHOULDER& RIGHT\_WRIST    &  $0.726$ \\
LEFT\_ELBOW    & LEFT\_WRIST     &  $0.722$ \\
RIGHT\_ELBOW   & RIGHT\_HIP      &  $0.721$ \\
LEFT\_HIP      & RIGHT\_ELBOW    & $-0.705$ \\
\bottomrule
\end{tabular}
\end{table}

In the XY plane, the dominant feature is a strong negative correlation
between left and right elbow ($r=-0.862$), alongside positive same-side
elbow--shoulder and ankle--knee coupling. In the XZ plane, the hip
anti-correlation ($r=-0.84$) and upper-limb coupling ($r=0.72$--$0.77$) are
notable. Figure~\ref{fig:correlation_heatmaps} visualises the full
correlation matrices.

\begin{figure*}[t]
  \centering
  \begin{subfigure}[b]{0.49\textwidth}
    \centering
    \includegraphics[width=\linewidth]{figs/fig06a_corr_theta.png}
    \caption{$\Delta\theta$ (XY plane)}
  \end{subfigure}\hfill
  \begin{subfigure}[b]{0.49\textwidth}
    \centering
    \includegraphics[width=\linewidth]{figs/fig06b_corr_psi.png}
    \caption{$\Delta\psi$ (XZ plane)}
  \end{subfigure}
  \caption{Correlation matrices of inter-joint direction angle errors.
    (a)~$\Delta\theta$ corresponds to Table~\ref{tab:correlation_theta};
    (b)~$\Delta\psi$ corresponds to Table~\ref{tab:correlation_psi}.
    Source: \texttt{06\_direction\_detection/output/correlation\_analysis/}.}
  \label{fig:correlation_heatmaps}
\end{figure*}

\subsection{Camera Viewpoint Dependency}

To characterise viewpoint dependence, we identified for each of the
107 frames the camera with the lowest and highest mean $|\Delta\theta|$
across joints (computed from
\texttt{11\_direction\_ditection/output/processed\_data/frame\_camera\_summary.csv},
107 frames $\times$ 201 cameras per frame).

The \emph{best} camera per frame achieves a mean $|\Delta\theta|$ of
$8.8\degree{}$ averaged over all frames (min $6.1\degree{}$, max $15.8\degree{}$), demonstrating
that near-optimal viewpoints exist. In contrast, the \emph{worst} camera
yields a mean of $80.1\degree{}$ (min $55.6\degree{}$, max $148.1\degree{}$), indicating that
highly unfavourable viewpoints produce errors comparable to random
direction estimates. This 9-fold difference between best and worst cameras
underscores the critical importance of camera placement in practical
deployments.

% ============================================================
% 5. DISCUSSION  (target: ~2 pages)
% ============================================================
\section{Discussion}

\subsection{Coordinate System Unification Effect}

The Y-axis orientation mismatch between Unity and MediaPipe causes a
systematic sign error in $\theta$ for every joint. Because
$\theta_{\text{MP}} \approx -\theta_{\text{GT}}$ when the Y-axis is
inverted, the uncorrected direction angle error satisfies
$\Delta\theta \approx -2\theta_{\text{GT}}$; for a joint at
$\theta_{\text{GT}} \approx 60\degree{}$ this yields $|\Delta\theta| \approx 120\degree{}$.
After applying $y \leftarrow -y$ to MediaPipe coordinates, the mean
$|\Delta\theta|$ for \emph{shoulder} joints dropped from $\approx120\degree{}$ to
$\approx10\degree{}$, and for \emph{elbow} joints from $\approx121\degree{}$ to
$\approx58\degree{}$~\cite{mediapipe_coord_issue}.

% NOTE: The "before Y-flip" values (~121\degree{} for elbow) are documented in
% 10_theta_verification/README.md and MEDIAPIPE_COORDINATE_SYSTEM.md.
% Precise per-joint before/after table requires re-running
% 10_theta_verification/test_04_coordinate_system.py without the Y-flip.

The residual $\approx58\degree{}$ elbow error after unification reflects genuine
estimation bias (training-data viewpoint imbalance, arm-at-side prior),
not coordinate system mismatch.

\subsection{Systematic Bias in Depth Estimation}

MediaPipe's BlazePose estimates each landmark via heatmaps and regression,
outputting 3-D world coordinates relative to the hip
centre~\cite{blazepose}. The architecture does not explicitly incorporate
anatomical constraints such as bone-length consistency or pelvic
rigidity. The mean $|\Delta\theta|$ for elbow joints after coordinate
unification remains $\approx58\degree{}$, indicating a persistent directional bias.

Many pose estimation models, including MediaPipe, are trained on public
datasets such as Human3.6M~\cite{h36m} and COCO~\cite{coco} in which
everyday, near-frontal, standing/walking poses dominate. The arm-at-side
prior embedded in these datasets likely biases the elbow direction estimate
toward a downward-hanging posture. The opposite-sign pattern in left vs.\
right elbow errors ($r=-0.862$) strongly suggests left--right flip data
augmentation during training; the standard deviation of elbow $\Delta\theta$
($\approx23$--$25\degree{}$) is small enough that the bias is relatively
predictable and correctable by post-processing.

\subsection{Left-Right Symmetric Hip Depth Error}

For hip joints, the mean $|\Delta\psi|\approx90\degree{}$ and the strong negative
correlation ($r=-0.8402$) indicate that the depth-direction errors are
anti-correlated: when one hip is estimated as forward, the other is
estimated as backward, creating a phantom pelvic-rotation artefact. The
identical absolute-value heatmap patterns for left and right hip
(Figure~\ref{fig:hip_psi_heatmap}) are consistent with
$\Delta\psi_{\text{L}} \approx -\Delta\psi_{\text{R}}$ at each observation.

In MediaPipe's skeleton model the pelvis is represented as the segment
connecting landmark~23 (left hip) and landmark~24 (right hip). This linear
representation without a multi-segment rigid structure (e.g., sacrum and
ilium) imposes no rigidity constraint on the relative depth of the two
hips during 3-D lifting. Consequently, even a small horizontal pixel
difference in the 2-D image can be interpreted as placing one hip closer
and the other farther from the camera, producing the symmetric error
pattern. Model improvements that impose pelvis rigidity as a hard or soft
constraint are needed.

\subsection{Upper-Limb Coupled Error}

The high positive correlations among shoulder, elbow, and wrist in
$\Delta\psi$ ($r=0.72$--$0.77$) suggest that errors propagate
hierarchically along the kinematic chain, or that the upper limb is
implicitly treated as a rigid unit during 3-D lifting. Incorporating
anatomical bone-length constraints via graph neural networks or explicit
kinematic models could reduce this propagation.

\subsection{Camera Viewpoint Dependency}

Selecting the best camera per frame reduces $|\Delta\theta|$ to 6.1\degree{}--15.8\degree{}
(mean 8.8\degree{}), while the worst camera yields 55.6\degree{}--148.1\degree{} (mean 80.1\degree{}).
This 9-fold range between best and worst viewpoints for the same body pose
confirms that estimation quality is heavily viewpoint-dependent.
The training data for BlazePose (primarily Human3.6M~\cite{h36m} and
COCO~\cite{coco}) is dominated by frontal and near-frontal viewpoints;
cameras at extreme angles (lateral, top-down) are underrepresented, leading
to reduced generalisation. This is consistent with the elevated MAE at
Y=2.0 (overhead cameras) observed in Table~\ref{tab:mae_by_layer}.

\subsection{Practical Implications}

MediaPipe provides useful 2-D joint localisation but its 3-D depth
estimation is fundamentally limited by monocular ambiguity.
For applications where relative joint angles in the image plane suffice
(e.g., stroke-count analysis in swimming, gross-motor rehabilitation
screening) MediaPipe is adequate. For applications requiring 3-D body
orientation (e.g., gait analysis, ergonomic risk assessment of torso
rotation) the systematic biases identified here---particularly in hip
$\Delta\psi$ and elbow $\Delta\theta$---must be compensated.

\subsection{Limitations}

This study has the following limitations. (1)~The avatar is a standard
Y-Bot mesh; skin deformation and clothing effects present in real subjects
are absent. (2)~Only a single walking motion was evaluated; other motion
types may produce different error patterns. (3)~The simulated lighting is
uniform; real-world illumination variation was not assessed. (4)~Only
MediaPipe was evaluated; comparison with other estimators (e.g., YOLO-Pose,
OpenPose) is left for future work.

% ============================================================
% 6. CONCLUSION  (target: ~0.5 page)
% ============================================================
\section{Conclusion}

We presented a large-scale multi-viewpoint evaluation of MediaPipe pose
estimation using a Unity simulation environment. The key findings are:
\begin{itemize}\itemsep0pt
  \item Joint angle MAE ranges from $\approx17\degree{}$ (knee) to $\approx40\degree{}$
        (shoulder); elbow and knee errors are the lowest.
        MAE increases at Y=2.0 (overhead) cameras, consistent with
        training-data viewpoint bias.
  \item Coordinate system unification (Y-axis flip) is essential:
        it reduces shoulder $|\Delta\theta|$ from $\approx120\degree{}$ to
        $\approx10\degree{}$ and elbow $|\Delta\theta|$ from $\approx121\degree{}$ to
        $\approx58\degree{}$.
  \item Elbow direction error ($|\Delta\theta|\approx58\degree{}$) reflects
        training-data bias and is systematic (SD $\approx23$--$25\degree{}$),
        offering scope for post-processing correction.
  \item Hip depth error ($|\Delta\psi|\approx90\degree{}$, $r=-0.84$) reflects
        an absence of pelvic rigidity constraints in the skeleton model.
  \item Upper-limb errors are strongly coupled ($r=0.72$--$0.77$),
        suggesting hierarchical error propagation.
\end{itemize}

Future directions include: (1)~incorporating pelvic rigidity constraints
into pose estimation models; (2)~fine-tuning on multi-viewpoint data
collected from diverse angles; (3)~extending the evaluation framework to
other estimators and motion types; and (4)~developing application-specific
correction filters for elbow and hip errors.

% ============================================================
\section*{Acknowledgments}
% ============================================================

This work was supported by the Japan Science and Technology Agency (JST)
under the Next-Generation Human Resource Development Program (Global Science
Campus Initiative).

% ============================================================
% REFERENCES
% ============================================================
\begin{thebibliography}{99}

% --- Original 3 references (kept from paper_en.md) ----------
\bibitem{mediapipe}
Google LLC, ``MediaPipe Pose Landmarker,''
\url{https://ai.google.dev/edge/mediapipe/solutions/vision/pose_landmarker},
2024.

\bibitem{mediapipe_coord_issue}
Google, ``Orientation of pose world landmarks,''
GitHub Issue \#3370, MediaPipe repository, 2022.
\url{https://github.com/google/mediapipe/issues/3370}

\bibitem{blazepose}
V.~Bazarevsky, I.~Grishchenko, K.~Raveendran, T.~Zhu, F.~Zhang,
and M.~Grundmann,
``BlazePose: On-device real-time body pose tracking,''
\textit{arXiv preprint arXiv:2006.10204}, 2020.

\bibitem{unity}
Unity Technologies, ``HumanBodyBones,''
\textit{Unity Scripting Reference}, 2024,
\url{https://docs.unity3d.com/ScriptReference/HumanBodyBones.html}.

\bibitem{h36m}
C.~Ionescu et al.,
``Human3.6M: Large Scale Datasets and Predictive Methods for 3D Human
Sensing in Natural Environments,''
\textit{IEEE Trans. Pattern Anal. Mach. Intell.}, vol.~36, no.~7,
pp.~1325--1339, 2014.

\bibitem{coco}
T.-Y.~Lin et al.,
``Microsoft COCO: Common Objects in Context,''
in \textit{Proc. ECCV}, 2014, pp.~740--755.

\bibitem{openpose}
Z.~Cao, G.~Hidalgo, T.~Simon, S.-E.~Wei, and Y.~Sheikh,
``OpenPose: Realtime Multi-Person 2D Pose Estimation Using Part Affinity
Fields,''
\textit{IEEE Trans. Pattern Anal. Mach. Intell.}, vol.~43, no.~1,
pp.~172--186, 2021.

% --- Independent BlazePose / MediaPipe evaluation ------------------
\bibitem{mroz2021}
S.~Mroz, N.~Baddour, C.~McGuirk, P.~Juneau, A.~Tu, K.~Cheung,
and E.~Lemaire,
``Comparing the quality of human pose estimation with BlazePose or
OpenPose,''
in \textit{Proc. 4th Int.\ Conf.\ Bio-Engineering for Smart Technologies
(BioSMART)}, 2021, pp.~1--4.
\url{https://doi.org/10.1109/BioSMART54244.2021.9677850}

% --- Synthetic / simulation-based GT datasets ----------------------
\bibitem{surreal}
G.~Varol, J.~Romero, X.~Martin, N.~Mahmood, M.~J.~Black, I.~Laptev,
and C.~Schmid,
``Learning from synthetic humans,''
in \textit{Proc. CVPR}, 2017, pp.~109--117.

\bibitem{agora}
P.~Patel, C.-H.~P.~Huang, J.~Tesch, D.~T.~Hoffmann, S.~Tripathi,
and M.~J.~Black,
``AGORA: Avatars in geography optimized for regression analysis,''
in \textit{Proc. CVPR}, 2021.

\bibitem{peoplesanspeople}
S.~E.~Ebadi et al.,
``PeopleSansPeople: A synthetic data generator for human-centric
computer vision,''
\textit{arXiv preprint arXiv:2112.09290}, 2021.

% --- Monocular 3D pose estimation / depth ambiguity ----------------
\bibitem{mehta2017}
D.~Mehta, H.~Rhodin, D.~Casas, P.~Fua, O.~Sotnychenko, W.~Xu,
and C.~Theobalt,
``Monocular 3D human pose estimation in the wild using improved CNN
supervision,''
in \textit{Proc. Int.\ Conf.\ 3D Vision (3DV)}, 2017, pp.~506--516.
\url{https://doi.org/10.1109/3DV.2017.00064}

\bibitem{vnect}
D.~Mehta et al.,
``VNect: Real-time 3D human pose estimation with a single RGB camera,''
\textit{ACM Trans.\ Graph.}, vol.~36, no.~4, pp.~44:1--44:14, 2017.

% --- Skeletal constraint / transformer-based models ----------------
\bibitem{poseformer}
C.~Zheng, S.~Zhu, M.~Mendieta, T.~Yang, C.~Chen, and Z.~Ding,
``3D human pose estimation with spatial and temporal transformers,''
in \textit{Proc. ICCV}, 2021, pp.~11656--11666.

\bibitem{motionbert}
W.~Zhu et al.,
``MotionBERT: A unified perspective on learning human motion
representations,''
in \textit{Proc. ICCV}, 2023.

\end{thebibliography}

\end{document}


